\begin{abstract}
General-purpose tokenizers may allocate vocabulary inefficiently for specialized
technical documentation, increasing sequence length. We present a controlled,
CPU-only case study of how training-corpus composition affects byte-level BPE
compression on unseen documentation. We hold the tokenizer algorithm,
pre-tokenizer, decoder, normalization, special tokens, vocabulary size, and
training-data budget fixed. Using byte-matched 32\,MiB corpora, we train
General-English, Azure-documentation, and 50/50 Mixed tokenizers with 8K, 16K,
and 32K vocabularies and evaluate them on nine leakage-controlled held-out
datasets. On an unseen cross-service Azure holdout containing 964 documents from
five wholly held-out services, the 16K Azure tokenizer uses \reduc{36.68} fewer
tokens per 1{,}000 UTF-8 bytes than the matched General tokenizer (95\%
paired-bootstrap CI $[35.77\%, 37.62\%]$; document win rate $1.000$). However, it
uses \reduc{34.75} more tokens on held-out general English. The 50/50 Mixed
tokenizer achieves a \reduc{34.21} reduction on the Azure holdout, retaining
\reduc{93.3} of the Azure-only tokenizer's in-domain gain, while using only
\reduc{4.39} more tokens on general English. The in-domain reduction is stable
across the evaluated vocabulary sizes, persists on prose-dominant content, and
transfers partially to other Microsoft and external Kubernetes documentation. We
evaluate tokenizer compression and CPU tokenization throughput only. We make no
claims about model quality, inference latency, GPU throughput, or realized cost.
We release code, pinned source manifests, per-document measurements, and a
machine-verified claim registry; raw source corpora are not redistributed.
\end{abstract}
